
\section{Koppler}
\renewcommand{\ort}{topics/Koppler}

\Term{Taktanregung}
\begin{pitemize}
    \item \Term{Gleichtaktanregung}: \cp
Signale gleicher Amplitude und gleicher Phase.
\begin{pitemize}
    \item $Z_e$
    \item $a_1^+ = \frac{1}{2} a_1,a_4^+ = \frac{1}{2} a_1$
\end{pitemize}
    \item \Term{Gegentaktanregung}: \cm
Signale gleicher Amplitude und engegengesetzer Phase.
\begin{pitemize}
    \item $Z_d$
    \item $a_1^- = \frac{1}{2} a_1,a_4^- = -\frac{1}{2} a_1$
\end{pitemize}
    \item Überlagerung: $a_1^+  + a_1^- = a_1$ , $a_4^+  + a_4^- = 0$
\end{pitemize}


\subsection{Leitungskoppler}

\bild{topics/Koppler/Leitungskoppler}

\bild{topics/Koppler/Bsp Leitungskoppler}

$|\Gamma|^2 = \frac{P2}{P1} \left\lvert\frac{S_{41}}{S_{32}S_{21}}\right\rvert^2$

\Term{Zweifach Symmetrischer Zweitor}:

\bild{topics/Koppler/Leitungskoppler 2}

Allseit. Anpassung: \\$S_{11}=S_{22}=S_{33}=S_{44}=0$\\
Symmetrien und Reziprok:\\
$S_{41}=S_{14}=S_{32}=S_{23}$\\
$S_{12}=S_{21}=S_{34}=S_{43}$\\

\Term{Mit Gleich- und Gegentaktanregung}:\\
$b_1= b_1^+ + b_1^- = S_{11}^+ a_1^+ + S_{11}^- a_1^- = \frac{1}{2} (S_{11}^+ + S_{11}^-) a_1$, $\rightarrow S_{11}$
$b_2= b_2^+ + b_2^- = S_{21}^+ a_1^+ + S_{21}^- a_1^- = \frac{1}{2} (S_{21}^+ + S_{21}^-) a_1$, $\rightarrow S_{21}$
$b_3= b_2^+ + b_2^- = S_{21}^+ a_4^+ + S_{21}^- a_4^- = \frac{1}{2} (S_{21}^+ - S_{21}^-) a_1$, $\rightarrow S_{31}$
$b_4= b_4^+ + b_4^- = S_{11}^+ a_4^+ + S_{11}^- a_4^- = \frac{1}{2} (S_{11}^+ - S_{11}^-) a_1$, $\rightarrow S_{41}$




\Term{Normierung der Widerstände}:
$z_e = \frac{Z_e}{Z_0}$,$z_d = \frac{Z_d}{Z_0}$ \\
$\begin{bmatrix}
    A^\ggt{+}{-} & B^\ggt{+}{-} \\ C^{\ggt{+}{-}} & D^{\ggt{+}{-}}
\end{bmatrix} = \begin{bmatrix}
    \cos(\beta l) & jz_{\ggt{e}{d}}\sin(\beta l) \\
    j\frac{1}{z_{\ggt{e}{d}}}\sin(\beta l) & \cos(\beta l)
\end{bmatrix}$

$S_{11}^+ = \frac{A^+ + B^+ - C^+ - D^+}{A^+ + B^+ + C^+ + D^+} = \frac{jz_e\sin(\beta l) -
    j\frac{1}{z_e}\sin(\beta l) }{2\cos(\beta l) + j(z_e + \frac{1}{z_e}) \sin(\beta l)}$\\

$S_{11}^- = \frac{j(z_d - \frac{1}{z_d})\sin(\beta l) }{2\cos(\beta l) + j(z_d + \frac{1}{z_d}) \sin(\beta l)}$
$S_{21}^+ = \frac{2}{A^+ + B^+ + C^+ + D^+} = \frac{2}{2\cos(\beta l) + j(z_e + \frac{1}{z_e}) \sin(\beta l)}$\\
$S_{21}^+ = \frac{2}{A^- + B^- + C^- + D^-} = \frac{2}{2\cos(\beta l) + j(z_d + \frac{1}{z_d}) \sin(\beta l)}$

Allseitige anpassung $S_{11} = 0 \rightarrow S_{11}^+ = - S_{11}^- \rightarrow z_e = \frac{1}{z_d} \leftrightarrow z_ez_d = 1 \rightarrow Z_eZ_d = Z_0^2$
Entkoppelte Tore: $S_{31} = \frac{1}{2}(S_{21}^+ - S_{21}^-) = 0 \rightarrow S_{21}^+ = S_{21}^-$
$S_{21} = S_{21}^+ = \frac{2}{2\cos(\beta l) + j(z_e + \frac{1}{z_e}) \sin(\beta l)}$\\
$S_{41} = \frac{1}{2}(S_{11}^+ - S_{11}^-) = S_{11}^+ = \frac{j(z_e - \frac{1}{z_e})\sin(\beta l) }{2\cos(\beta l) + j(z_e + \frac{1}{z_e}) \sin(\beta l)}$

$\angle(S_{21}, S_{41}) = 90^\circ$\\
Mehrstufig wenn höhere Breitband.

\subsection{Ringkoppler}

\bild{topics/Koppler/Ringkoppler}

\begin{pitemize}
    \item 3dB-Kopplung: 1 $\rightarrow$ 2 und 1 $\rightarrow$ 3
    \item Entkopplung: 1 und 4
    \item Einspeisung an 1
    \item $Z_1 = ?$, $Z_2 = ?$ 
\end{pitemize}


\begin{pitemize}
    \item \Term{Gleichtakt}
    \begin{pitemize}
        \item \cp  an 1 und 4: $a_1^+ = \frac{1}{2}a_1, a_4^+ = \frac{1}{2}a_1$ 
        \item \bild{topics/Koppler/Ringkoppler Gleichtakt}
    \end{pitemize}
    \item \Term{Gegentakt}
    \begin{pitemize}
        \item \cm  an 1 und 4: $a_1^- = \frac{1}{2}a_1, a_4^- = -\frac{1}{2}a_1$
        \item \bild{topics/Koppler/Ringkoppler Gegentakt}
    \end{pitemize}
\end{pitemize}


\bild{topics/Koppler/Leitungsdarstellung}
\begin{pitemize}
    \item $Z_{ein} = Z_L \frac{Z_A + jZ_L \tan(\beta l)}{Z_L + jZ_A \tan(\beta l)}$
    \item $\tan(\beta l) = \tan(\frac{2\pi}{\lambda} \frac{\lambda}{8}) = \tan(\frac{\pi}{4}) = 1$
    \item \cp $Z_A \rightarrow \infty \Rightarrow Z_{ein,LL}^+ = \lim_{Z_A \rightarrow \infty}\limits Z_{ein} = -jZ_L = -jZ_1$
    \item \cm $Z_A \rightarrow 0 \Rightarrow Z_{ein,LL}^- = jZ_L = jZ_1 = j\frac{1}{Y_1}$
\end{pitemize}

Kettenmatrizen fortsetzung

$\left[M^\ggt{+}{-}\right] = \begin{bmatrix}
    1 & 0 \\ -jY_1Z_0 & 1
\end{bmatrix}\begin{bmatrix}
    0 & j\frac{Z_2}{Z_0} \\ j\frac{Z_0}{Z_2} & 0
\end{bmatrix}\begin{bmatrix}
    1 & 0 \\ -jY_1Z_0 & 1
\end{bmatrix}
= \begin{bmatrix}
    \ggt{+}{-}Z_2Y_1 & j \frac{Z_2}{Z_0}\\ j\frac{Z_0}{Z_2} - jY_1^2Z_0Z_2 & \ggt{+}{-}Y_1Z_2
\end{bmatrix} = \begin{bmatrix}
    A^\ggt{+}{-} & B^\ggt{+}{-} \\ C^\ggt{+}{-} & D^\ggt{+}{-}
\end{bmatrix}$

Anforderung: Allseitige Anpassung
\begin{pitemize}
    \item $S_{11} = \frac{1}{2}(S_{11}^+ + S_{11}^-) \overset{!}{=} 0$
    \item $S_{41} = S_{14} = \frac{1}{2}(S_{11}^+ - S_{11}^-) \overset{!}{=} 0$
    \item $S_{11}^+ = \frac{A^+ + B^+ - C^+ - D^+}{A^+ + B^+ + C^+ + D^+} \overset{!}{=} 0$
    \item Da $A^+ = D^+ \Rightarrow B^+ = C^+$
    \item $\frac{Z_2}{Z_0} = \frac{Z_0}{Z_2} - Y_1^2 Z_2 Z_0$
    \item $|S_{21}| = \lvert\frac{1}{2} \left(S_{21}^+ + S_{21}^-\right)\rvert = |S_{31}| = \lvert\frac{1}{2} \left(S_{21}^+ - S_{21}^-\right)\rvert$
    \item $\frac{Z_2}{Z_0} = Y_1 Z_2 \Rightarrow Y_1 = \frac{1}{Z_0}$
    \item $1 = \frac{Z_0^2}{Z_2^2} - Y_1^2Z_0^2 \Rightarrow Z_2 = \frac{1}{\sqrt{2}}Z_0$
\end{pitemize}

\begin{pitemize}
    \item $\angle(S_{21}, S_{31}) = 90^\circ$
    \item 3dB-Kopplung: 1 $\rightarrow$ 2 und 1 $\rightarrow$ 3
    \item Entkopplung: 1 und 4
\end{pitemize}


\subsubsection{Rat-Race-Koppler}
\bild{topics/Koppler/Rat-Race-Koppler}
\begin{pitemize}
    \item Einsp. 1: 3 entkoppelt, $\angle(S_{21}, S_{41}) = 0^\circ$
    \item Einsp. 3: 1 entkoppelt, $\angle(S_{23}, S_{43}) = 180^\circ$
\end{pitemize}








\subsection{Resistiver Koppler}
\bild{topics/Koppler/ResKop}

Synthese: All. Anpassung
\begin{pitemize}
    \item $S_{31} = S_{13} = S_{24} = S_{42} = 0$
\end{pitemize}
$S_{11} = \frac{Z_{ein} -Z_0}{Z_{ein} + Z_0}$

\cp und \cm an 1 und 4.\\

\bildf{r\_koppler}

\cp $\Rightarrow$ in der Symmetrieeben $\Rightarrow$ LL
\begin{pitemize}
    \item \bildf{r\_koppler\_plus}
    \item $S_{11}^+ = \frac{2R_2 + Z_0 - Z_0}{2 R_2 + Z_0 + Z_0} = \frac{R_2}{R_2+Z_0}$
\end{pitemize}

\cm
\begin{pitemize}
    \item \bildf{r\_koppler\_minus}
    \item $S_{11}^- = -\frac{Z_0}{R_1 + Z_0}$
\end{pitemize}

$S_{11} = \frac{1}{2}(S_{11}^+ + S_{11}^-) \overset{!}{=} 0 \Rightarrow S_{11}^+ = -S_{11}^-$ \\
$\frac{R_2}{R_2 + Z_0} = \frac{Z_0}{R_1 + Z_0} \Rightarrow R_1R_2 = Z_0^2$

$S_{41} = \frac{1}{2}(S_{11}^--S_{11}^-) = S_{11}^+ = \frac{R_2}{R_2 + Z_0}$\\
$S_{21} = \frac{Z_0}{Z_0 + R_2}$


Häufige Wahl: $R_1 = R_2 = Z_0 \Rightarrow S_{21}=S_{41} = \frac{1}{2}$

Nullgradkoppler: $\angle(S_{21}, S_{41}) = 0^\circ$

\subsection{Transformatorkoppler}

\bildf{trafo\_koppler}

\begin{pitemize}
    \item Transformator sei ideal
    \item aus Symmetrie sind 1 und 4 entkoppelt
    \item bei Einspeisung an 1 $\Rightarrow$ Signale an 2 und 3 in Gegenphase
    \item bei Einspeisung an 4 $\Rightarrow$ Signale an 2 und 3 in Phase
    \item $\Rightarrow$ Signale an 2 und 3 entkoppelt mit $Z = \frac{Z_0}{2}$
\end{pitemize}


\subsection{Nullgradkoppler/Signalteiler}

\bildf{NullGradKoppler}
\bildf{null\_grad\_koppler\_EB}
\begin{pitemize}
    \item \cp und \cm and 2 und 3
    \item \cp an 2 und 3:
    \begin{pitemize}
        \item \bildf{null\_grad\_koppler\_plus}
        \item $Z_{ein} \frac{Z_1}{2Z_0} \overset{!}{=} 0 \Rightarrow Z_1 = \sqrt{2}Z_0$
    \end{pitemize}
    \item \cm an 2 und 3:
    \begin{pitemize}
        \item \bildf{null\_grad\_koppler\_minus}
        \item $\frac{R}{2}=Z_0 \Rightarrow R = 2Z_0$
    \end{pitemize}
\end{pitemize}